%% lncs-switch.tex -- machinery for the LNCS (Eurocrypt) build of the paper.
%%
%% Loaded by main.tex after the packages and before macros.tex.  The switch
%% itself, \iflncs, is set at the very top of main.tex (before \documentclass).
%%
%% Markup available in main.tex.  Every command below is a no-op in the full
%% version, so the full version is exactly the paper as before.
%%
%%   \begin{lncsappendix}[Heading]{name}
%%      ...
%%   \end{lncsappendix}
%%      LNCS: the block is removed from where it stands and re-inserted,
%%      verbatim, in the supplementary material at \lncsappendixspine, in
%%      order of appearance.  The optional [Heading] emits \section{Heading}
%%      in front of the block (for blocks that do not start with their own
%%      \section).  The \begin and \end lines must stand alone on their line.
%%
%%   \begin{lncsstatement}{Name}
%%      \begin{theorem}[Title]\label{t:...} ... \end{theorem}
%%   \end{lncsstatement}
%%      LNCS: the theorem is stated (and numbered) in the body, at the point
%%      where \lncsplace{Name} stands, and restated -- same heading, same
%%      number (taken from the first \label inside the block), "(restated)"
%%      added, labels suppressed -- where it originally stood, provided that
%%      place is inside an lncsappendix block.  Full version: typeset in place.
%%
%%   \begin{lncstable}{Name} ... \end{lncstable}
%%      Same as lncsstatement for tables/figures/any block: LNCS moves it to
%%      \lncsplace{Name} and leaves nothing behind.
%%
%%   \lncsplace[setup]{Name}   LNCS: typeset the block Name here (inside a
%%                         group, after running the optional setup code, e.g.
%%                         \renewcommand's the block relies on).  Full: nothing.
%%   \lncsappendixspine    LNCS: emit all lncsappendix blocks (put it after
%%                         \appendix).  Full: nothing.
%%   \iflncs ... \else ... \fi      for anything else (front matter etc.).
%%
%% How it works (LNCS build only): before \begin{document} the source file is
%% read once, line by line, and every marked block is copied verbatim into
%% \jobname-lncs-app-<name>.tex / \jobname-lncs-stmt-<Name>.tex (regenerated
%% at every run; safe to delete).  In the main pass the marked blocks are
%% skipped in place (line based, like the comment package) and the generated
%% files are \input where requested.  Nothing is captured as macro arguments,
%% so verbatim material, comment environments and catcode changes inside the
%% blocks are safe.

\makeatletter

%% ---------------------------------------------------------------------------
%% 1. llncs compatibility with the paper's theorem set-up (amsthm + aliascnt in
%%    macros.tex).  llncs predefines theorem-like environments and counters;
%%    we drop the ones macros.tex redefines, so that it can define them as
%%    usual (the others -- case, exercise, note, problem, property, solution --
%%    stay as llncs defines them and are unused).
%% ---------------------------------------------------------------------------
\iflncs
  \@for\lncs@e:=theorem,lemma,proposition,corollary,claim,conjecture,question,definition,convention,construction,remark,fact,example\do{%
    \expandafter\let\csname\lncs@e\endcsname\relax
    \expandafter\let\csname end\lncs@e\endcsname\relax
    \@removefromreset{\lncs@e}{chapter}%
    \@removefromreset{\lncs@e}{section}%
    \expandafter\let\csname c@\lncs@e\endcsname\relax
    \expandafter\let\csname the\lncs@e\endcsname\relax
  }%
  \def\@thmcountersep{.}% llncs empties it; the paper numbers theorems as <section>.<n>
\fi

%% ---------------------------------------------------------------------------
%% 2. Collaborator notes (\albert, \albertimportant, \john, ...) are hidden in
%%    the LNCS build unless \lncsshownotestrue is set in main.tex.
%% ---------------------------------------------------------------------------
\newif\iflncsshownotes
\AtBeginDocument{%
  \iflncs\iflncsshownotes\else
    \@for\lncs@e:=albert,albertimportant,albertforkai,john,ulrich,remco,Albert,kai,xander,psiv,osdnk,Aru\do{%
      \ifcsname\lncs@e\endcsname\expandafter\renewcommand\csname\lncs@e\endcsname[1]{}\fi
    }%
  \fi\fi
}

%% ---------------------------------------------------------------------------
%% 3. The block markup.
%% ---------------------------------------------------------------------------
\iflncs
  %% 3a. In-place removal of the blocks, reusing comment.sty's line-based
  %%     exclusion (\excludecomment); we only add the argument(s) of \begin.
  \excludecomment{lncsappendix}
  \excludecomment{lncsstatement}
  \excludecomment{lncstable}
  \expandafter\let\expandafter\lncs@skip@appendix\csname lncsappendix\endcsname
  \expandafter\let\expandafter\lncs@skip@statement\csname lncsstatement\endcsname
  \expandafter\let\expandafter\lncs@skip@table\csname lncstable\endcsname
  \expandafter\def\csname lncsappendix\endcsname{\@ifnextchar[{\lncs@skip@appendix@i}{\lncs@skip@appendix@i[]}}
  \def\lncs@skip@appendix@i[#1]#2{\lncs@skip@appendix}
  \expandafter\def\csname lncsstatement\endcsname#1{\lncs@skip@statement}
  \expandafter\def\csname lncstable\endcsname#1{\lncs@skip@table}

  %% 3b. Pre-scan of the source: copies the blocks to files (expl3).
  \ExplSyntaxOn
  \ior_new:N  \l__lncs_src_ior
  \iow_new:N  \l__lncs_app_iow
  \iow_new:N  \l__lncs_blk_iow
  \seq_new:N  \l__lncs_m_seq
  \seq_new:N  \g__lncs_regions_seq
  \seq_new:N  \g__lncs_blocks_seq
  \prop_new:N \g__lncs_labels_prop
  \tl_new:N   \l__lncs_label_tl
  \bool_new:N \l__lncs_in_region_bool
  \bool_new:N \l__lncs_in_block_bool
  \str_new:N  \l__lncs_block_kind_str
  \str_new:N  \l__lncs_block_name_str
  \str_new:N  \l__lncs_src_str
  \int_new:N  \l__lncs_line_int

  \msg_new:nnn {lncs} {no-source} {LNCS~pre-scan:~cannot~open~the~source~file~'#1'.}
  \msg_new:nnn {lncs} {unclosed}  {LNCS~pre-scan:~a~'#1'~block~was~never~closed.}
  \msg_new:nnn {lncs} {nested}    {LNCS~pre-scan,~line~#2:~'#1'~cannot~be~nested.}
  \msg_new:nnn {lncs} {stray-end} {LNCS~pre-scan,~line~#2:~\iow_char:N\\end{#1}~without~\iow_char:N\\begin.}
  \msg_new:nnn {lncs} {duplicate} {LNCS~pre-scan,~line~#2:~the~name~'#1'~is~used~twice.}
  \msg_new:nnn {lncs} {no-block}  {No~lncsstatement/lncstable~block~named~'#1'~exists~in~the~source.}

  \cs_new_protected:Npn \lncs_prescan:n #1
    {
      \str_set:Nn \l__lncs_src_str {#1}
      \ior_open:NnTF \l__lncs_src_ior {#1}
        {
          \bool_set_false:N \l__lncs_in_region_bool
          \bool_set_false:N \l__lncs_in_block_bool
          \int_zero:N \l__lncs_line_int
          \ior_str_map_inline:Nn \l__lncs_src_ior
            {
              \int_incr:N \l__lncs_line_int
              \__lncs_line:n {##1}
            }
          \ior_close:N \l__lncs_src_ior
          \bool_if:NT \l__lncs_in_block_bool  { \msg_error:nnn {lncs} {unclosed} { \l__lncs_block_kind_str } }
          \bool_if:NT \l__lncs_in_region_bool { \msg_error:nnn {lncs} {unclosed} {lncsappendix} }
        }
        { \msg_error:nnn {lncs} {no-source} {#1} }
    }

  % One source line (#1 is a string: catcode-12 tokens, spaces).
  \cs_new_protected:Npn \__lncs_line:n #1
    {
      \regex_extract_once:nnNTF
        { \A \\begin\{lncsappendix\} (?: \[ ([^\]]*) \] )? \{ ([^\}]*) \} \s* \Z } {#1} \l__lncs_m_seq
        { \__lncs_region_begin:xx { \seq_item:Nn \l__lncs_m_seq {2} } { \seq_item:Nn \l__lncs_m_seq {3} } }
        {
          \regex_match:nnTF { \A \\end\{lncsappendix\} \s* \Z } {#1}
            { \__lncs_region_end: }
            {
              \regex_extract_once:nnNTF
                { \A \\begin\{(lncsstatement|lncstable)\} \{ ([^\}]*) \} \s* \Z } {#1} \l__lncs_m_seq
                { \__lncs_block_begin:xx { \seq_item:Nn \l__lncs_m_seq {2} } { \seq_item:Nn \l__lncs_m_seq {3} } }
                {
                  \regex_match:nnTF { \A \\end\{(lncsstatement|lncstable)\} \s* \Z } {#1}
                    { \__lncs_block_end: }
                    { \__lncs_content:n {#1} }
                }
            }
        }
    }

  \cs_new_protected:Npn \__lncs_content:n #1
    {
      \bool_if:NTF \l__lncs_in_block_bool
        {
          \iow_now:Nn \l__lncs_blk_iow {#1}
          \prop_if_in:NVF \g__lncs_labels_prop \l__lncs_block_name_str
            {
              \regex_extract_once:nnNT { \\label\{([^\}]*)\} } {#1} \l__lncs_m_seq
                { \prop_gput:NVx \g__lncs_labels_prop \l__lncs_block_name_str { \seq_item:Nn \l__lncs_m_seq {2} } }
            }
        }
        { \bool_if:NT \l__lncs_in_region_bool { \iow_now:Nn \l__lncs_app_iow {#1} } }
    }

  \cs_new_protected:Npn \__lncs_region_begin:nn #1#2   % #1 heading (may be empty), #2 name
    {
      \bool_if:NT \l__lncs_in_region_bool { \msg_error:nnx {lncs} {nested} {lncsappendix} { \int_use:N \l__lncs_line_int } }
      \seq_if_in:NnT \g__lncs_regions_seq {#2} { \msg_error:nnx {lncs} {duplicate} {#2} { \int_use:N \l__lncs_line_int } }
      \bool_set_true:N \l__lncs_in_region_bool
      \seq_gput_right:Nn \g__lncs_regions_seq {#2}
      \iow_open:Nn \l__lncs_app_iow { \c_sys_jobname_str -lncs-app- #2 .tex }
      \iow_now:Nx \l__lncs_app_iow
        { \c_percent_str ~ lncsappendix~block~'#2',~copied~from~\l__lncs_src_str \c_space_tl line~ \int_use:N \l__lncs_line_int ~ (generated~file,~do~not~edit) }
      \tl_if_empty:nF {#1}
        { \iow_now:Nx \l__lncs_app_iow { \c_backslash_str section \c_left_brace_str #1 \c_right_brace_str } }
    }
  \cs_generate_variant:Nn \__lncs_region_begin:nn { xx }

  \cs_new_protected:Npn \__lncs_region_end:
    {
      \bool_if:NF \l__lncs_in_region_bool { \msg_error:nnx {lncs} {stray-end} {lncsappendix} { \int_use:N \l__lncs_line_int } }
      \bool_if:NT \l__lncs_in_block_bool   { \msg_error:nnn {lncs} {unclosed} { \l__lncs_block_kind_str } }
      \bool_set_false:N \l__lncs_in_region_bool
      \iow_close:N \l__lncs_app_iow
    }

  \cs_new_protected:Npn \__lncs_block_begin:nn #1#2   % #1 kind, #2 name
    {
      \bool_if:NT \l__lncs_in_block_bool { \msg_error:nnx {lncs} {nested} {#1} { \int_use:N \l__lncs_line_int } }
      \seq_if_in:NnT \g__lncs_blocks_seq {#2} { \msg_error:nnx {lncs} {duplicate} {#2} { \int_use:N \l__lncs_line_int } }
      \bool_set_true:N \l__lncs_in_block_bool
      \str_set:Nn \l__lncs_block_kind_str {#1}
      \str_set:Nn \l__lncs_block_name_str {#2}
      \seq_gput_right:Nn \g__lncs_blocks_seq {#2}
      \iow_open:Nn \l__lncs_blk_iow { \c_sys_jobname_str -lncs-stmt- #2 .tex }
      \iow_now:Nx \l__lncs_blk_iow
        { \c_percent_str ~ #1~block~'#2',~copied~from~\l__lncs_src_str \c_space_tl line~ \int_use:N \l__lncs_line_int ~ (generated~file,~do~not~edit) }
    }
  \cs_generate_variant:Nn \__lncs_block_begin:nn { xx }

  \cs_new_protected:Npn \__lncs_block_end:
    {
      \bool_if:NF \l__lncs_in_block_bool { \msg_error:nnx {lncs} {stray-end} {lncsstatement/lncstable} { \int_use:N \l__lncs_line_int } }
      \bool_set_false:N \l__lncs_in_block_bool
      \iow_close:N \l__lncs_blk_iow
      \bool_if:NT \l__lncs_in_region_bool
        {
          \str_if_eq:VnT \l__lncs_block_kind_str {lncsstatement}
            { \iow_now:Nx \l__lncs_app_iow { \c_backslash_str lncsrestate \c_left_brace_str \l__lncs_block_name_str \c_right_brace_str } }
        }
    }

  % User level.
  \cs_generate_variant:Nn \seq_if_in:NnTF { Nx }
  \cs_new_protected:Npn \lncs_place:n #1
    {
      \seq_if_in:NxTF \g__lncs_blocks_seq { \tl_to_str:n {#1} }
        { \exp_args:Nx \file_input:n { \c_sys_jobname_str -lncs-stmt- #1 .tex } }
        { \msg_error:nnn {lncs} {no-block} {#1} }
    }
  \cs_new_protected:Npn \lncs_restate:n #1
    {
      \seq_if_in:NxTF \g__lncs_blocks_seq { \tl_to_str:n {#1} }
        {
          \group_begin:
            \prop_get:NnNF \g__lncs_labels_prop {#1} \l__lncs_label_tl { \tl_clear:N \l__lncs_label_tl }
            \tl_set_eq:NN \lncs@restated@label \l__lncs_label_tl
            \cs_set_eq:NN \label \use_none:n
            \lncs@restated@envs
            \exp_args:Nx \file_input:n { \c_sys_jobname_str -lncs-stmt- #1 .tex }
          \group_end:
        }
        { \msg_error:nnn {lncs} {no-block} {#1} }
    }
  \cs_new_protected:Npn \lncs_spine:
    {
      \seq_map_inline:Nn \g__lncs_regions_seq
        { \exp_args:Nx \file_input:n { \c_sys_jobname_str -lncs-app- ##1 .tex } }
    }
  \NewDocumentCommand \lncsplace { O{} m } { \group_begin: #1 \lncs_place:n {#2} \group_end: }
  \NewDocumentCommand \lncsrestate { m } { \lncs_restate:n {#1} }
  \NewDocumentCommand \lncsappendixspine { } { \lncs_spine: }

  % Run the pre-scan now.  \lncssourcefile may be predefined (e.g. by a wrapper
  % file with a different \jobname); it defaults to main.tex.
  \providecommand\lncssourcefile{main.tex}
  \exp_args:Nx \lncs_prescan:n { \lncssourcefile }
  \ExplSyntaxOff

  %% 3c. Restatements: unnumbered amsthm environments whose heading repeats the
  %%     original heading and number.  Inside \lncsrestate the theorem-like
  %%     environments of macros.tex are redirected to these.
  \def\lncs@restated@label{}
  \def\lncs@restated@head#1{#1\ifx\lncs@restated@label\@empty\else~\ref{\lncs@restated@label}\fi}
  \theoremstyle{plain}
  \newtheorem*{lncs@restated@plain}{\lncs@restated@head\lncs@restated@name}
  \theoremstyle{definition}
  \newtheorem*{lncs@restated@definition}{\lncs@restated@head\lncs@restated@name}
  \theoremstyle{plain}
  \def\lncs@restated@note#1{\ifx\@empty#1\@empty restated\else#1, restated\fi}
  \def\lncs@redirect#1#2#3{% env name, printed name, style (skipped if the env does not exist)
    \ifcsname #1\endcsname
      \expandafter\renewenvironment\expandafter{#1}[1][]{\def\lncs@restated@name{#2}\begin{lncs@restated@#3}[\lncs@restated@note{##1}]}{\end{lncs@restated@#3}}%
    \fi}
  \def\lncs@restated@envs{%
    \lncs@redirect{theorem}{Theorem}{plain}%
    \lncs@redirect{lemma}{Lemma}{plain}%
    \lncs@redirect{proposition}{Proposition}{plain}%
    \lncs@redirect{corollary}{Corollary}{plain}%
    \lncs@redirect{claim}{Claim}{plain}%
    \lncs@redirect{conjecture}{Conjecture}{plain}%
    \lncs@redirect{question}{Question}{plain}%
    \lncs@redirect{fact}{Fact}{plain}%
    \lncs@redirect{definition}{Definition}{definition}%
    \lncs@redirect{convention}{Convention}{definition}%
    \lncs@redirect{construction}{Construction}{definition}%
    \lncs@redirect{remark}{Remark}{definition}%
    \lncs@redirect{example}{Example}{definition}%
  }
\else
  %% Full version: the markup does nothing.  The environments deliberately do
  %% not open a group, so definitions inside them behave as if unmarked.
  \newcommand\lncsappendix{\@ifnextchar[{\lncs@noop@i}{\lncs@noop@i[]}}
  \def\lncs@noop@i[#1]#2{\endgroup}
  \def\endlncsappendix{\begingroup\def\@currenvir{lncsappendix}}
  \newcommand\lncsstatement[1]{\endgroup}
  \def\endlncsstatement{\begingroup\def\@currenvir{lncsstatement}}
  \newcommand\lncstable[1]{\endgroup}
  \def\endlncstable{\begingroup\def\@currenvir{lncstable}}
  \newcommand\lncsplace[2][]{}
  \newcommand\lncsrestate[1]{}
  \newcommand\lncsappendixspine{}
\fi

\makeatother
